\subsection{Del catálogo de realizaciones al inventario de 471 masas}
La extensión universal sigue cinco operaciones:

\[
\begin{aligned}
 \text{identidad física}
 &\longrightarrow
 \text{proyector de representación}
 \longrightarrow
 \text{registro genealógico elemental o compuesto}\\
 &\longrightarrow
 \text{operador/polo}
 \longrightarrow
 \text{carta de comparación}.
\end{aligned}
\]
Para estados elementales, el segundo mapa actúa sobre una fibra del atlas. Para mesones y bariones, actúa sobre las 432 y 1728 expresiones compuestas. Para resonancias, el cuarto mapa devuelve un polo complejo. Para sectores topológicos, devuelve primero un brecha espectral y una representación de trenza. Esta tipificación permite abarcar todo el inventario sin declarar que cada nombre externo corresponde a una celda elemental.

\Needspace{7\baselineskip}
\subsection{De las 384 anchuras al operador de supervivencia}
Cada anchura requiere

\[
 X\longmapsto Q_XUQ_X
 \longmapsto \rho(Q_XUQ_X)
 \longmapsto\Gamma_X
\]
o, equivalentemente, el polo de \(H_{\rm eff}\). El lector debe incorporar canales de decaimiento, multiplicidad, umbral y espacio de fases. La igualdad

\[
 \Gamma_{\bar X}=\Gamma_X
\]
se deduce sólo si el operador de decaimiento es covariante bajo conjugación.

\Needspace{7\baselineskip}
\section{Masa, temperatura y estadística}
\Needspace{7\baselineskip}
\subsection{Transductor térmico entre carácter de masa, temperatura y ocupación estadística}
La constante de Boltzmann es el transductor

\[
 k_B:\mathcal L_\Theta\longrightarrow\mathcal L_E.
\]
El reloj CT--108 fija la identidad

\[
 k_B\Theta_{\rm clk}\log3
 =\hbar_{\rm int}\frac{2\pi}{108t_0}.
\]
Esta ecuación vincula una unidad de información trítica, una frecuencia y una energía. No introduce la coordenada SI de \(k_B\) en la ley; determina el producto tipado una vez elegida la carta térmica.

\Needspace{7\baselineskip}
\subsection{Proyección fermiónica del atlas y distribución de Fermi–Dirac}
Para niveles

\[
 \epsilon=(0,3,6,9),\qquad
 g=(1,6,12,8),
\]
y restricciones

\[
 (N,E)=(12,66),
\]
la completación exterior produce

\[
 2\,104\,494
\]
microestados fermiónicos compatibles. Los multiplicadores efectivos son

\[
 (\beta,\mu)
 =(0.1530701135,\;4.5028651976).
\]
Esta realización ejemplifica cómo masa, energía y exclusión emergen de rutas fermiónicas sin que la estadística sea un ajuste adicional.

\Needspace{7\baselineskip}
\subsection{Proyección bosónica del atlas y distribución de Bose–Einstein}
Sobre los mismos niveles y restricciones, la completación simétrica produce

\[
 259\,436\,430
\]
microestados, con

\[
 (\beta,\mu)
 =(0.0545672796,\;-15.9256978191).
\]
La diferencia entre Fermi--Dirac y Bose--Einstein es, por tanto, una diferencia en la regla de composición de rutas, no en el valor de una masa individual.

\Needspace{7\baselineskip}
\subsection{Criterio de condensación de las secciones bosónicas y ocupación macroscópica}
Un condensado bosónico aparece cuando una componente simétrica adquiere ocupación macroscópica en el autovalor mínimo del operador. Un mar fermiónico, en cambio, llena un conjunto de autovalores distintos hasta una frontera. HMT describe ambos mediante el mismo espectro de masa/energía y dos funtores de Fock distintos. Esta separación impide identificar “muchas partículas en una ruta” con una partícula más masiva.

\Needspace{7\baselineskip}
\section{Propagación diferencial de incertidumbres y estabilidad estructural del carácter de masa}
\Needspace{7\baselineskip}
\subsection{Invariancia del grafo generativo bajo el contraste metrológico posterior}
La salida interna se congela en el orden

\[
 \mathcal R_{324}
 \to P_X
 \to\mathcal L_X
 \to\sigma_X
 \to\eta_X
 \to\widehat M_X
 \to\mu_X.
\]
Sólo después se abre la carta externa. Una permutación arbitraria de las 471 masas o de las 384 anchuras debe dejar invariantes todos los objetos anteriores.

\Needspace{7\baselineskip}
\subsection{Descomposición de la incertidumbre estructural en contribuciones experimental, de carta, ruta, torre y acoplamiento}
La incertidumbre total se descompone como

\[
 \Sigma_{\rm tot}
 =\Sigma_{\rm ext}
 +\Sigma_{\rm carta}
 +\Sigma_{\rm route}
 +\Sigma_{\rm tower}
 +\Sigma_{\rm coupling}.
\]
Aquí:

\begin{itemize}
\item \(\Sigma_{\rm ext}\) es la covarianza experimental;
\item \(\Sigma_{\rm carta}\) cuantifica esquema, escala y conversión de unidades;
\item \(\Sigma_{\rm route}\) compara todas las rutas de la órbita admisible;
\item \(\Sigma_{\rm tower}\) usa la cota de cola de la serie armónica;
\item \(\Sigma_{\rm coupling}\) mide la variación entre entrelazadores permitidos.
\end{itemize}
No se suman necesariamente como números independientes; la ecuación expresa una descomposición de fuentes que debe convertirse en una covarianza conjunta.

\Needspace{7\baselineskip}
\subsection{Derivada del observable de masa respecto de la firma y de los lectores}
Para

\[
 \log m_X=
 \log m_\star+
 \log\chi_0(\sigma_X)
 +\kappa_X\log R_{\rm act}
 +\sum_h\eta_{X,h}s_h,
\]
la variación es

\[
 d\log m_X
 =d\log m_\star
 +\langle\Theta_0,d\sigma_X\rangle
 +\langle\sigma_X,d\Theta_0\rangle
 +\log R_{\rm act}\,d\kappa_X
 +\kappa_X\,d\log R_{\rm act}
 +\sum_h(s_h\,d\eta_{X,h}+\eta_{X,h}\,ds_h).
\]
Esta fórmula permite atribuir cada decimal a ruta, constante, lector, torre o carta. Un resultado es robusto cuando perturbaciones que preservan las premisas mantienen su clase espectral y su incertidumbre calculada.

\Needspace{7\baselineskip}
\subsection{Ensayos de dependencia estructural mediante supresión controlada de hoja, orientación, acarreo, memoria, devanado y frontera}
Una prueba de dependencia elimina por separado hoja, orientación, acarreo, memoria, devanado o frontera y vuelve a calcular el objeto. El resultado esperado no es que todo cambie, sino que cambien exactamente las coordenadas dependientes de la pieza eliminada. Si suprimir memoria deja invariante un coeficiente definido por retorno, la construcción de ese coeficiente queda refutada.

\Needspace{7\baselineskip}
\subsection{Obstrucciones a la identificación de rutas, caracteres y magnitudes metrológicas}
La ley queda refutada en una aplicación concreta si ocurre cualquiera de estos hechos:

\begin{enumerate}
\item una masa externa cambia la ruta o el proyector elegido;
\item la firma publicada no se regenera desde el registro genealógico;
\item el refinamiento no converge o depende de una truncación elegida por el residual;
\item dos rutas equivalentes reciben masas distintas sin ruptura de simetría;
\item una antipartícula adquiere masa distinta sin violación declarada de conjugación;
\item un triplete de color se escinde sin ruptura de color;
\item se publica un escalar aunque la varianza espectral sea positiva;
\item una anchura se obtiene del carácter de masa sin operador temporal;
\item una carta quark omite esquema o escala;
\item una inestabilidad se interpreta como taquión persistente.
\end{enumerate}
\Needspace{7\baselineskip}
